% ============================================================
% Section: Introduction
% ============================================================
\section{Introduction}
\label{sec:introduction}

% --- Layer 1: The Challenge ---
Nuclear reactor design and safety analysis depend on accurate neutron flux predictions.
The neutron diffusion equation---a tractable approximation to the Boltzmann transport equation---remains the standard model for core calculations, and its two-group variant, separating fast and thermal neutrons, retains enough physics to capture moderation and absorption while admitting analytical treatment in simple geometries~\cite{duderstadt1976nuclear,bell1970nuclear,lee2020nuclear}.

Solutions to the multigroup diffusion equations have traditionally come from two directions.
Analytical methods---Fourier series, Green's functions, eigenvalue decomposition---yield exact results for canonical geometries and remain the primary verification benchmarks~\cite{zanette2018fourier,williams2013fourier}, but they are restricted to regular domains with simple material layouts; deriving them for coupled multigroup systems under non-trivial boundary conditions requires considerable manual effort.
Numerical discretizations (FDM, FEM) handle arbitrary geometries and heterogeneous media at the price of discretization error and, for full-core problems, substantial computation~\cite{lee2020nuclear,kuridan2023neutron}.

% --- Layer 1 continued: ML paradigm ---
Scientific machine learning has opened a third route.
Physics-Informed Neural Networks (PINNs) fold PDE residuals into the training loss, producing mesh-free solutions for both forward and inverse problems~\cite{raissi2019pinn}.
In reactor physics, successive architectural changes have reduced PINN errors on two-group problems~\cite{elhareef2022pinn,tlpinn2023,naspinn2025}.
Neural operator approaches, including DeepONet~\cite{lu2021deeponet} and the Fourier Neural Operator~\cite{li2021fno}, learn solution operators that generalize across problem parameters, achieving large speedups over traditional solvers in neutron transport surrogate modeling~\cite{kobayashi2024deeponet}.
Both families have limits: PINNs exhibit training instability and accuracy bottlenecks (0.6--15\% error for two-group problems~\cite{elhareef2022pinn}), and neural operators are constrained by the coverage and fidelity of their training data.

% --- Layer 2: What We Do ---
This paper combines classical analytical derivation with LLM-driven program evolution to solve the two-group neutron diffusion equation on a two-dimensional finite homogeneous domain with boundary sources.
The first component is a Fourier cosine series solution that gives exact ground truth for the boundary-source problem in rectangular geometry.
The second is \textbf{FaMoU} (\underline{Fa}st \underline{Mo}de m\underline{U}tation), an evolutionary program synthesis loop---inspired by FunSearch and AlphaEvolve---in which a large language model iteratively generates and refines solver programs, scored against the analytical ground truth.
We benchmark both components alongside FDM, high-order FDM, and PINN on the same problem instance.

% --- Layer 3: Technical Approach ---
Specifically, we derive the two-group neutron flux distribution by expanding the solution in a Fourier cosine series that satisfies the homogeneous Neumann boundary conditions, with the boundary source incorporated through Fourier decomposition of the inhomogeneous condition.
The resulting coupled system of ordinary differential equations is solved analytically for each Fourier mode, yielding pointwise ground truth with controllable accuracy governed by the number of retained modes.
For the evolutionary component, inspired by FunSearch~\cite{funsearch2024} and AlphaEvolve~\cite{alphaevolve2025}, an initial seed program is iteratively improved by a large language model, with each candidate scored by an automated evaluator that measures PDE and boundary-condition residuals as RMS values at fixed test points.
A distinguishing feature of our setup is that the evaluator computes error against an \emph{exact} analytical ground truth, rather than relying on heuristic scores or combinatorial bounds as in most program evolution work.

% --- Layer 4: Contributions ---
Our main contributions are as follows:
\begin{itemize}
    \item \textbf{Analytical benchmark.} We derive and implement the complete analytical solution for the two-group neutron diffusion equation with boundary sources in a 2D finite homogeneous rectangular medium, providing a high-precision reference solution that satisfies the PDE exactly in the continuous limit and achieves residuals below $4 \times 10^{-5}$ under numerical evaluation.

    \item \textbf{First LLM-driven evolution for physics PDE solving.} To our knowledge, this is the first work to apply the LLM-driven program evolution paradigm---pioneered by FunSearch~\cite{funsearch2024} and AlphaEvolve~\cite{alphaevolve2025} for mathematical and algorithmic discovery---to the domain of computational physics PDE solving. We demonstrate that evolutionary program synthesis can discover solver programs achieving competitive accuracy with hand-crafted numerical methods.

    \item \textbf{Comprehensive method comparison.} We present a unified benchmark comparing analytical, FDM, PINN, and LLM-evolved solvers on the same problem instance, providing quantitative insights into the accuracy--efficiency trade-offs across four distinct methodological paradigms.

    \item \textbf{Overflow-safe adaptive truncation.} We identify and resolve a numerical overflow issue in the standard Fourier implementation, deriving a stability criterion that allows the expansion to include all float64-safe modes (up to $n{=}221$), improving the combined score from $0.9989$ to $0.9997$.
\end{itemize}

The remainder of this paper is organized as follows.
\Cref{sec:related_work} surveys related work across classical numerical methods, PINNs, neural operators, and LLM-driven program evolution.
\Cref{sec:problem} presents the mathematical formulation and analytical solution.
\Cref{sec:method} describes the LLM-driven evolutionary framework and baseline methods.
\Cref{sec:experiments} reports experimental results and comparative analysis.
\Cref{sec:conclusion} concludes with a discussion of implications and future directions.
